Zur genauen Bestimmung der Erdmasse benützt man heutzutage die Daten von Satellitenbahnen. Die Tabelle zeigt einige Beispiele:

\begin{center}
\begin{tabular}{lrrr}
\hline
Satellit & \multicolumn{2}{c}{Abstand zur Erdoberfläche} & Umlaufzeit \\
 & im Perigäum & im Apogäum & \\
\hline
Nimbus & \hpNiO & \haNiO & \TNiO \\
Skynet 2 & \hpSkO & \haSkO & \TSkO \\
Nato 1 & \hpNaO & \haNaO & \TNaO \\
\hline
\end{tabular}
\end{center}

Berechnen Sie aus diesen Angaben zunächst den mittleren Abstand (entspricht der grossen Halbachse) der Satelliten von der Erde und daraus die Erdmasse.
